\documentclass[11pt,letterpaper]{article}

\usepackage[utf8]{inputenc}
\usepackage[spanish]{babel}
\usepackage[margin=2.5cm]{geometry}
\usepackage{amsmath,amssymb}
\usepackage{enumitem}
\usepackage{titlesec}
\usepackage{array}
\usepackage{hyperref}
\usepackage{xcolor}

\definecolor{unicaucaverde}{RGB}{0,90,60}
\hypersetup{colorlinks=true,linkcolor=unicaucaverde,urlcolor=unicaucaverde}

\titleformat{\section}{\normalfont\Large\bfseries\color{unicaucaverde}}{\thesection}{1em}{}
\titleformat{\subsection}{\normalfont\large\bfseries}{\thesubsection}{1em}{}

\setlength{\parindent}{0pt}
\setlength{\parskip}{6pt}

\title{\textcolor{unicaucaverde}{\Large Guía de Aprendizaje --- Semana 1}\\[4pt]
  \large Sistemas embebidos en FPGA: del sistema Nios II en Quartus al software en Nios II EDS}
\author{Énfasis --- Universidad del Cauca}
\date{2026}

\begin{document}
\maketitle

\begin{center}
\begin{tabular}{|>{\bfseries}p{4cm}|p{10cm}|}
\hline
Semana & 1 de 16 \\ \hline
Unidad & Introducción a sistemas embebidos en FPGA con procesador soft-core Nios II \\ \hline
Subtemas & Arquitectura Nios II y flujo de herramientas (Quartus Prime, Platform Designer,
  Nios II EDS); construcción y prueba de un sistema Nios II mínimo sobre la tarjeta DE2-115 \\ \hline
Horas de clase & 4 \\ \hline
Resultado de aprendizaje & RA1 (parte 1 de 2; continúa en Semana 2 con periféricos memory-mapped
  propios y E/S sobre PIO) \\ \hline
\end{tabular}
\end{center}

\section*{Relevancia para ingeniería}

Un \emph{soft processor} como Nios II es un procesador completo descrito en HDL que se sintetiza
dentro del tejido lógico de una FPGA, en vez de venir fabricado como bloque fijo (\emph{hard core})
en el silicio. Eso permite diseñar sistemas donde el procesador convive, en el mismo chip, con
lógica digital hecha a la medida --- controladores propios, aceleradores, interfaces no
estándar --- comunicados por el mismo bus. Es el punto de entrada natural al co-diseño
hardware/software: la habilidad de decidir qué parte de un sistema conviene resolver en software
sobre un procesador y qué parte conviene resolver en lógica dedicada, y de mover esa frontera según
el desempeño, el área o la potencia que exija el proyecto. Esa misma habilidad es la que se necesita
para trabajar con SoCs FPGA comerciales (Cyclone V SoC, Zynq) y, más adelante, con el sucesor de
Nios II en Intel/Altera --- Nios V, basado en RISC-V.

\section{Objetivo de la semana}

Construir, compilar y ejecutar un sistema mínimo basado en el procesador soft-core Nios II sobre la
FPGA Cyclone IV E de la tarjeta DE2-115, recorriendo el flujo completo de herramientas ---
Quartus Prime $\to$ Platform Designer $\to$ Nios II EDS --- y verificar su funcionamiento con un
programa \texttt{Hello World} en C que se comunica con el computador anfitrión por la JTAG UART.

\section{Resultados de aprendizaje de la semana}

\begin{itemize}[leftmargin=1.2cm]
  \item[\textbf{RA1.1.}] Explicar la arquitectura de un procesador soft-core Nios II (núcleos
        \texttt{/e}, \texttt{/s}, \texttt{/f}, bus Avalon-MM) y el rol de cada herramienta del
        flujo de diseño --- Quartus Prime, Platform Designer, Nios II EDS --- en el paso de
        ``sistema descrito en hardware'' a ``programa en software ejecutándose sobre ese
        hardware''.
  \item[\textbf{RA1.2.}] Crear un proyecto en Quartus Prime, ensamblar en Platform Designer un
        sistema Nios II mínimo (núcleo, memoria on-chip, JTAG UART), integrarlo en un diseño
        top-level, asignar pines para la DE2-115, compilar y programar la FPGA, y luego compilar y
        ejecutar sobre ese sistema un programa en C usando Nios II EDS.
\end{itemize}

\section{Contenidos}

\begin{itemize}[leftmargin=1.2cm]
  \item \textbf{1.1.} Procesadores soft-core vs. hard-core; arquitectura Nios II y bus Avalon-MM.
  \item \textbf{1.2.} El flujo de herramientas: Quartus Prime, Platform Designer, Nios II EDS.
  \item \textbf{1.3.} Del sistema en hardware al software: BSP y compilación cruzada.
  \item \textbf{1.4.} Práctica integradora: sistema Nios II mínimo y \texttt{Hello World} sobre la
        DE2-115.
\end{itemize}

\section{Plan de la sesión (4 horas)}

\begin{enumerate}[leftmargin=1.2cm]
  \item[\textbf{Hora 1.}] Procesadores soft-core vs. hard-core; por qué tiene sentido sintetizar un
        procesador dentro de una FPGA. Arquitectura Nios II (núcleos \texttt{/e}, \texttt{/s},
        \texttt{/f}) y el bus Avalon Memory-Mapped. Panorama de la tarjeta DE2-115 (FPGA Cyclone IV
        E EP4CE115F29C7, oscilador de 50\,MHz, pulsadores, LEDs, USB-Blaster integrado) y del flujo
        Quartus Prime $\to$ Platform Designer $\to$ Nios II EDS de extremo a extremo (teoría +
        demostración en vivo).
  \item[\textbf{Hora 2.}] Taller guiado: crear el proyecto en Quartus Prime, abrir Platform
        Designer y ensamblar el sistema --- núcleo Nios II, fuente de reloj, memoria on-chip, JTAG
        UART, System ID --- conectando los buses Avalon y asignando direcciones base
        automáticamente; generar el sistema (HDL + \texttt{.sopcinfo}).
  \item[\textbf{Hora 3.}] Integrar el sistema generado en el diseño top-level de Quartus (instanciar
        el módulo generado, conectar reloj y reset), importar la asignación oficial de pines de la
        DE2-115, compilar el proyecto completo (Analysis \& Synthesis, Fitter, Assembler) y
        programar la FPGA por USB-Blaster.
  \item[\textbf{Hora 4.}] Taller: en Nios II EDS, crear el proyecto de aplicación y BSP a partir de
        la plantilla \texttt{Hello World} usando el \texttt{.sopcinfo} generado, compilar, ejecutar
        sobre la tarjeta y verificar la salida por la terminal Nios II (JTAG UART); resolución de
        fallas comunes (reloj/reset sin conectar, BSP desincronizado con el hardware, tiempo de
        espera de programación agotado).
\end{enumerate}

\section{Actividades}

\begin{itemize}[leftmargin=1.2cm]
  \item \textbf{En clase}: talleres de las horas 2, 3 y 4 (ver arriba).
  \item \textbf{Trabajo autónomo}: agregar un componente PIO de salida al sistema Qsys, regenerar el
        BSP y modificar \texttt{hello\_world.c} para hacer parpadear un LED de la DE2-115 además de
        imprimir el mensaje; dibujar el diagrama de bloques del sistema armado (núcleo, memoria,
        JTAG UART, PIO, y las conexiones Avalon-MM entre ellos). Entrega antes de la Semana 2.
  \item \textbf{Software/hardware}: Quartus Prime 18.1 Standard Edition (última versión con soporte
        para Cyclone IV y para el flujo clásico Nios II/Nios II EDS), tarjeta Terasic DE2-115 y
        cable USB-Blaster.
\end{itemize}

\section{Material de apoyo}

\begin{itemize}[leftmargin=1.2cm]
  \item \texttt{notas.tex} --- desarrollo teórico completo de 1.1 a 1.4, con el procedimiento
        detallado del laboratorio.
  \item \texttt{diapositivas.tex} --- síntesis visual de la sesión.
\end{itemize}

\section{Evaluación de la semana}

Entregable funcional al inicio de la Semana 2: captura de pantalla (o video corto) de la terminal
Nios II mostrando el mensaje de \texttt{Hello World} recibido desde la tarjeta, más el diagrama de
bloques del sistema Qsys armado, para verificar RA1.1 y RA1.2 antes de avanzar a periféricos propios
sobre PIO.

\section{Bibliografía de la semana}

\begin{enumerate}[leftmargin=1.2cm]
  \item Intel/Altera. \emph{Nios II Processor Reference Guide}.
  \item Intel/Altera. \emph{Nios II Software Developer's Handbook}.
  \item Intel/Altera. \emph{Platform Designer (Qsys) System Design Tutorial}.
  \item Terasic Technologies. \emph{DE2-115 User Manual}.
\end{enumerate}

\end{document}
